\begin{table}[t]
    \centering
    \begin{tabular}{@{}lcccc@{}}
        \toprule
        \textbf{Method} & \textbf{fraction defined} & \textbf{coverage} & \textbf{mean bound} & \textbf{mean truth} \\
        \midrule
        \fixedcp (pre-registered budget) & 1.00 & 0.626 & 0.697 & 0.682 \\
        \adaptivecp (data-chosen stop)   & 1.00 & 0.638 & 0.676 & 0.682 \\
        \unioncp (anytime-valid)         & 1.00 & \textbf{0.671} & 0.726 & 0.682 \\
        \extrapolative (power law)       & 0.89 & 0.555 & 0.686 & 0.682 \\
        \bottomrule
    \end{tabular}
    \caption{\textbf{E3: nulls are badly calibrated, and sequential rigour does not repair them.}
    Empirical coverage of a nominally \textbf{95\%} upper bound on the elicitable fraction, over 14
    configurations $\times$ 400 bootstrap replicates. All four methods under-cover badly. Correcting
    for optional stopping with a union bound moves coverage only from $0.638$ to $0.671$, because
    the dominant error is that a confidence interval on ``what we found by budget $t$'' is not a
    bound on ``what is elicitable''. Closest to nominal in \textbf{bold}.}
    \label{tab:e3_coverage}
\end{table}
